% ==================== 摘要 ====================
{\bfseries 摘要}\quad
以单晶硅（$a=b=c=5.43$~\AA，$\alpha=\beta=\gamma=90^\circ$）为例，在 Origin 中建立了一套可复用的晶体学计算模板：由晶格参数计算晶胞体积、倒易点阵参数、晶面间距、倒易矢量长度及晶面夹角，并据此构建了 Cu K$\alpha$ 条件下的单晶硅 PDF 卡片。把带轴电子衍射模拟分为两部分：先由 VESTA 导出的原子坐标和结构因子计算晶面的消光系数，判断哪些反射允许、哪些消光；再以 [001]、[111]、[332]、[334]、[345] 五个带轴为例，用晶带方程建立二维倒易点阵，计算衍射点坐标，并用散点图和三维 Waterfall 图绘制衍射图。结果表明，模拟图谱的对称性与点阵几何符合立方金刚石结构的预期，110 等反射因结构因子为零而消光。

\vspace{4pt}

\par\noindent{\heiti 关键词：}晶体学计算；Origin；PDF 卡片；电子衍射；单晶硅；结构因子
